\subsection{信号预处理与网络训练}
\label{sec:experiment-training}

应用于所有采集数据的信号预处理流程遵循第\ref{sec:preprocessing}节所述的协议。原始激光超声信号以有限次平均（$N = 16$）采集，经FIR带通滤波器（100\,kHz--2\,MHz）处理以保留兰姆波频段，随后采用矩形时窗选择以预测$S_0$模态到达时刻为中心的$30\,\mu\text{s}$窗口。对窗口化段落应用逐发$z$-分数归一化，以标准化跨采集会话的振幅。所有预处理后的信号被重采样至固定长度$L = 1024$以匹配PMDF-Net的输入维度。时频分支的STFT参数（汉宁窗$W = 256$，跳步$H = 64$，75\%重叠）在训练和推理阶段保持不变。

\textbf{配对数据集构建。}训练数据按照第\ref{sec:dataset}节的协议以配对格式组装。输入信号由$N = 16$次平均的激光超声采集组成，目标信号由相同空间位置处$M = 256$次平均的采集组成。输入与目标之间$10\log_{10}(M/N) \approx 12$\,dB的信噪比差异建立了网络学习的去噪目标。

\textbf{数据增强。}训练期间，每个配对样本经过随机化的逐样本增强以提高泛化能力：随机时间偏移量为信号长度的$\pm 12.5\%$，以相对于干净信号[-6, +6]\,dB范围内的信噪比水平注入加性高斯噪声，振幅缩放均匀采样于$[0.8, 1.2]$。增强独立应用于有限平均输入和高度平均目标以保持配对完整性。验证和测试期间不应用增强。

\textbf{训练配置。}PMDF-Net使用AdamW优化器~\cite{Loshchilov2019}训练，初始学习率为$1 \times 10^{-3}$，解耦权重衰减为$1 \times 10^{-4}$，小批量大小为$32$。学习率在完整训练过程中遵循从$1 \times 10^{-3}$至$1 \times 10^{-4}$的余弦退火计划。融合层应用概率$p = 0.2$的Dropout。当验证重建损失在连续$15$个epoch内无改善时训练终止，恢复经过验证的最佳模型权重。模型包含约$4.2 \times 10^6$个可训练参数。在单块NVIDIA RTX 4090（24\,GB显存）上，完整的180个epoch训练运行约需4.5小时，每个epoch平均吞吐量约为$2.3 \times 10^4$个训练步。